\section{Discussion}\label{sec:discussion}

The three limits of Section~\ref{sec:three-limits} translate into specific
information for the decisions that a solver makes about OBBT. We summarize
what the results of this paper support, and what they do not.

\paragraph{Starting.} A round cannot move the endpoints of a box that is
already fixed for the given relaxation and cutoff, although its dual
solutions can still be useful, for example for Lagrangian variable
bounds~\cite{gleixner2017-three-enhancements-for-optimization-based}. A
finite certificate proves fixedness (Section~\ref{sec:certificates}) and
remains valid at every larger cutoff. In a box that strictly contains the
certified box, it no longer proves a stall, but it still bounds the total
movement of all later rounds. Two structural situations explain when fixed
boxes are present. If a model has several global minimizers, for example by
symmetry, the iterates contain their hull and cannot converge to a point
(Lemma~\ref{lem:sublevel-hull}), although they can still remove parts of
the box outside the hull. If on every face of a box shape some point has
negative tangent-model value, its second-order relaxation gap exceeds the
objective's second-order growth there. Boxes of that shape around a minimizer
are then fixed at every sufficiently small scale and every cutoff $U\ge f^*$
(Theorem~\ref{thm:local-stall}); for quadratic objectives with McCormick
relaxations, condition~(a) of Proposition~\ref{prop:quadratic-rows} gives
this at every scale. That condition concerns the relaxation, not the
problem: the strongly convex objectives of Example~\ref{ex:many-term} stall
under termwise McCormick relaxations but not under a relaxation that keeps
the convex quadratic whole. Other results guarantee that a round moves the
box: the box threshold of Proposition~\ref{prop:box-threshold}, a stored
dual bound below the current endpoint
(Proposition~\ref{prop:dual-envelope-con} and
Corollary~\ref{cor:cutoff-multiplier-con}), and the contraction conditions
of Section~\ref{sec:local-rates}. None of them guarantees that the movement
shortens the search, and Example~\ref{ex:history} shows that the history of
earlier calls cannot identify a first action that is guaranteed to be best.
The precursor study of Section~\ref{sec:experiments} illustrates the cost
of starting blindly: more than nine tenths of its first-round time went to
the 109 instances whose nonlinear variables are all integer, and on only 25
of them did the first round change a bound.

\paragraph{Stopping.} A protected box inside the current box bounds the total
remaining movement of every endpoint under all later rounds of the
specified operator at cutoffs not below its threshold. When it equals the
current box, the remaining movement is zero and stopping that operator
loses nothing it could have achieved. The local theory shows when such
certificates can exist. In the stall regime they exist at every
sufficiently small scale. In the strictly contracting regime with a
positive cutoff slack and the closedness property~\eqref{eq:closed-family},
the limit is a fixed box of width $O(\sqrt\epsilon)$, so a certificate
exists at the limit, and the iterates reach its neighborhood after a number
of rounds that grows only logarithmically in $1/\epsilon$
(Section~\ref{sec:local-rates-consequences}). With zero slack, no box of
positive width is fixed near the minimizer. An iterate that reaches the
minimizer can be certified as a fixed singleton. When the optimal incumbent
lies in the current box, its protected singleton
(Corollary~\ref{cor:original-points}) bounds the remaining movement, and the
bound is exact if the iterates converge to that point
(Example~\ref{ex:sharp-halving}). Without such an inner box, a residual
certificate (Section~\ref{sec:residual}) can justify a tolerance stop.
These regimes need not cover every relaxation family
(Example~\ref{ex:critical-scalar}). Small observed changes certify nothing
beyond the bound that nesting already gives: the ratio of successive widths
can tend to one during slow contraction, near a positive floor, or at a
stall, and valid monotone families with identical histories can have
different futures (Section~\ref{sec:local-rates-consequences} and
Proposition~\ref{prop:history}). The exception is an exact complete round
that changes nothing, which proves that the box is fixed. Stopping rules based
on observed progress are therefore heuristics, which is how the measured
policy uses them.

\paragraph{Repeating after a change.} Under the contraction hypotheses, the
width of the limit box is at most of order $\sqrt\epsilon$ and the
remaining relaxation gap at most of order $\epsilon$
(Corollaries~\ref{cor:tangent-rate} and~\ref{cor:local-gap}); both orders
are attained in the two-variable quadratic model
(Proposition~\ref{prop:two-variable-cutoff-floor}). An improved incumbent is
therefore the main reason to reconsider OBBT at a box where it has already
converged: it can shrink the first limit, the hull of the original sublevel
set, and it can remove fixed boxes of the second. A certified stop remains
valid while the cutoff stays at or above the face threshold of its pool.
Below that threshold the stored proof expires, which is a precise event at
which reconsidering becomes justified, though not a proof that it will pay.
The certified box itself can remain fixed below the face threshold, down to
the box threshold of Proposition~\ref{prop:box-threshold}, whose computation
requires at most $2n$ face optimizations. For a single frozen round, retained relaxation
points and stored dual vectors bracket the possible and the guaranteed
change of each support after a cutoff decrease
(Section~\ref{sec:constraints}). A branching step keeps a certificate valid
in every child that contains the certified box and requires a new check in a
child that does not. Prior work already reevaluates stored dual
inequalities when the incumbent improves; the certificates add a statement
about the operator's remaining movement, not a new trigger mechanism.

\paragraph{Width, certificates, and search cost.} The guarantees of this
paper bound the movement of a specified tightening operator, or the change
of its relaxation bound. They do not bound search time, and the evidence
shows that the difference matters. In the prospective comparison, a policy
that solved $20.9\%$ fewer auxiliary LPs spent more time in its propagator.
The adaptive arm recorded more calls, each with its own construction and
checking cost; neither added policy produced an additional solve or a
demonstrated speedup (Section~\ref{sec:experiments}). In the precursor study,
tightened boxes reduced node counts, with the largest mean reductions in
groups whose mean control time was a few seconds, while the savings in the
final solve did not offset the OBBT cost. Both findings
agree with Section~\ref{sec:effort}: a budget on optional work bounds that
work, but no budget statement connects the enhanced search with an unchanged
baseline. The value of a stopping certificate is therefore to avoid work that
cannot change the box. Whether that saves time depends on how much the
certificate costs. Some certificates cost little beyond an exact check, such
as the singleton of an exactly feasible incumbent in the current box or the
optima of a completed exact round when they pass the rebuilt-row check. If finding a certificate costs about as
much as a round, reuse across rounds, cutoffs, or nodes may be needed to
recover that cost. The measured policy used no certificates of either
kind, and the experiments do not assess them. On the
present evidence, native SCIP remains the appropriate default and the added
policy is an experimental implementation.

\paragraph{Scope.} Every rate and stall statement belongs to a relaxation and
a formulation. At zero cutoff slack, the convex quadratic $\sum_ix_i^2$ with
$n\ge2$ stalls when each square is relaxed by the McCormick inequalities of
the product $x_i\cdot x_i$, contracts by the factor $1/2$ per round when the
valid rows $y_i\ge0$ are added, and collapses in one round under its exact
convex relaxation (Example~\ref{ex:finite-square}); eliminating a linear
equation before relaxing changes the rate in the example of
Section~\ref{sec:constraints}. The results on iterated tightening, rates,
and future-round certificates assume monotonicity of the whole relaxation
construction. A pool of cuts separated at earlier solutions does not
guarantee it, and neither do tangent points that move with the box, as in
the measured policy (Section~\ref{sec:foundations}). The current-round
ceilings, the validated dual bounds, and the results for one fixed LP do
not need it. The local theory involves the minimizer and tangent data
and is not by itself an activation rule. The exact reference procedure
certifies its own rational relaxation family; the numerical policy validates
each applied bound with directed rounding, but its cutoff comes from a
numerically accepted incumbent, and a complete SCIP solve is not an exact
certificate. The prospective comparison uses a small, selected cohort with
short time limits on a shared machine and supports only the conclusions
stated in Section~\ref{sec:experiments}.

Three extensions would change the practical picture. Certificates found
cheaply from points that a native solver produces anyway could make the
all-future guarantees usable at low cost. Tangent models for lifted
relaxations with auxiliary variables in the constraints would extend the
local theory to the relaxations that solvers build. A validated model of how
reductions of width or of the relaxation gap change the later search would be
needed before any certificate of remaining benefit can guide a choice among
tightening, cutting, and branching.
